\documentclass[12pt,a4paper]{article}

% -------------------- Packages --------------------
\usepackage[margin=1in]{geometry}
\usepackage{amsmath}
\usepackage{setspace}
\usepackage{times}
\usepackage{titlesec}
\usepackage{booktabs}
\usepackage{array}
\usepackage{longtable}
\usepackage{hyperref}
\usepackage{xcolor}
\usepackage{listings}

% -------------------- Formatting --------------------
\onehalfspacing
\setlength{\parindent}{0pt}
\setlength{\parskip}{6pt}

\titleformat{\section}
  {\normalfont\Large\bfseries}
  {\thesection.}{0.5em}{}

\titleformat{\subsection}
  {\normalfont\large\bfseries}
  {\thesubsection.}{0.5em}{}

\hypersetup{
    colorlinks=true,
    linkcolor=black,
    urlcolor=blue,
    citecolor=black
}

\lstdefinestyle{jsonstyle}{
    basicstyle=\ttfamily\small,
    breaklines=true,
    frame=single,
    backgroundcolor=\color{gray!5},
    showstringspaces=false
}

\lstset{style=jsonstyle}

% -------------------- Document --------------------
\begin{document}

% -------------------- Title Page --------------------
\begin{titlepage}
    \centering

    \vspace*{1.5cm}

    {\Large \textbf{Galala University}}\\[1.5cm]

    {\LARGE \textbf{Civara}}\\[0.4cm]

    {\Large \textbf{API Contract}}\\[0.5cm]

    {\large \textbf{A Public Statistical Analysis API for Student Performance Datasets}}\\[2cm]

    \begin{table}[h!]
        \centering
        \renewcommand{\arraystretch}{1.4}
        \begin{tabular}{>{\bfseries}l l}
            Student Name: & Mahros Mohamed \\
            Student ID: & 223106831 \\
            Course: & Probability and Statistics \\
            University: & Galala University \\
            Project Name: & Civara \\
            Domain: & \texttt{civara.mahros.dev} \\
        \end{tabular}
    \end{table}

    \vfill

    {\large \today}

\end{titlepage}

% -------------------- Main Content --------------------

\section{Overview}

Civara is a simple public API designed to analyze student performance datasets. The system allows a user to upload a CSV file, sends the dataset to the API, and receives a statistical analysis response.

The API focuses on descriptive statistics, confidence intervals, and hypothesis testing for comparing two independent groups, such as male and female students.

The API does not require authentication and does not depend on a database. Analysis is performed per request, and any persistence may be handled temporarily through application memory or browser cache.

\section{Base URL}

The production base URL is:

\begin{lstlisting}
https://civara.mahros.dev/api/v1
\end{lstlisting}

For local development, the base URL may be:

\begin{lstlisting}
http://localhost:3030/api/v1
\end{lstlisting}

\section{Authentication}

The API is public.

No authentication token, API key, or user account is required.

\section{Data Persistence}

The API does not permanently store uploaded datasets.

The expected behavior is:

\begin{itemize}
    \item The client uploads a dataset.
    \item The server analyzes the dataset.
    \item The server returns the analysis result immediately.
    \item The client may store the result in browser cache or local storage.
    \item The server may keep temporary in-memory data only when needed.
\end{itemize}

\section{Supported File Format}

The API accepts CSV files.

The uploaded CSV file should contain the following columns:

\begin{table}[h!]
\centering
\renewcommand{\arraystretch}{1.4}
\begin{tabular}{p{4cm} p{8cm}}
\toprule
\textbf{Column Name} & \textbf{Description} \\
\midrule
\texttt{gender} & The group variable used for comparison. \\
\texttt{math score} & Student score in mathematics. \\
\texttt{reading score} & Student score in reading. \\
\texttt{writing score} & Student score in writing. \\
\bottomrule
\end{tabular}
\caption{Required CSV columns}
\end{table}

The API calculates the average score using:

\[
\text{Average Score} =
\frac{\text{Math Score} + \text{Reading Score} + \text{Writing Score}}{3}
\]

\section{API Endpoints}

\section{Standard API Response Envelope}

All endpoints in this API return a consistent JSON envelope.

\subsection{Success Response (Single Object)}

\begin{lstlisting}
{
  "success": true,
  "message": "Request completed successfully",
  "data": {
    // ...
  },
  "meta": null
}
\end{lstlisting}

\subsection{Success Response (List + Pagination)}

When an endpoint returns a list, it uses the same envelope, and includes pagination metadata.

\begin{lstlisting}
{
  "success": true,
  "message": "Request completed successfully",
  "data": [
    // ...
  ],
  "meta": {
    "pagination": {
      "page": 1,
      "limit": 10,
      "totalItems": 125,
      "totalPages": 13,
      "hasNextPage": true,
      "hasPreviousPage": false
    }
  }
}
\end{lstlisting}

\subsection{Error Response}

All error responses also follow the same envelope and include an \texttt{error} object.

\begin{lstlisting}
{
  "success": false,
  "message": "Validation failed",
  "data": null,
  "meta": null,
  "error": {
    "code": "VALIDATION_ERROR",
    "message": "The uploaded file is missing the required column: gender."
  }
}
\end{lstlisting}

\subsection{Health Check}

\textbf{Endpoint}

\begin{lstlisting}
GET /health
\end{lstlisting}

\textbf{Description}

Checks whether the API is running.

\textbf{Successful Response}

\begin{lstlisting}
{
  "success": true,
  "message": "Request completed successfully",
  "data": {
    "status": "ok",
    "service": "Civara API"
  },
  "meta": null
}
\end{lstlisting}

\textbf{Status Codes}

\begin{table}[h!]
\centering
\renewcommand{\arraystretch}{1.3}
\begin{tabular}{c p{9cm}}
\toprule
\textbf{Status Code} & \textbf{Meaning} \\
\midrule
200 & API is running successfully. \\
\bottomrule
\end{tabular}
\caption{Health check status code}
\end{table}

\newpage

\subsection{Analyze Dataset}

\textbf{Endpoint}

\begin{lstlisting}
POST /analyze
\end{lstlisting}

\textbf{Description}

Uploads a CSV dataset and returns statistical analysis results.

\textbf{Request Type}

\begin{lstlisting}
multipart/form-data
\end{lstlisting}

\textbf{Request Body}

\begin{table}[h!]
\centering
\renewcommand{\arraystretch}{1.4}
\begin{tabular}{p{4cm} p{3cm} p{6cm}}
\toprule
\textbf{Field} & \textbf{Type} & \textbf{Description} \\
\midrule
\texttt{file} & File & The CSV dataset to be analyzed. \\
\texttt{scoreType} & String & Optional. Defaults to \texttt{average}. \\
\bottomrule
\end{tabular}
\caption{Analyze dataset request body}
\end{table}

\textbf{Supported Score Types}

\begin{table}[h!]
\centering
\renewcommand{\arraystretch}{1.4}
\begin{tabular}{p{4cm} p{8cm}}
\toprule
\textbf{Value} & \textbf{Description} \\
\midrule
\texttt{average} & Uses the average of math, reading, and writing scores. \\
\texttt{math} & Uses math score only. \\
\texttt{reading} & Uses reading score only. \\
\texttt{writing} & Uses writing score only. \\
\bottomrule
\end{tabular}
\caption{Supported score types}
\end{table}

\section{Successful Analysis Response}

\textbf{Status Code}

\begin{lstlisting}
200 OK
\end{lstlisting}

\textbf{Response Body}

\begin{lstlisting}
{
  "success": true,
  "message": "Request completed successfully",
  "data": {
    "project": "Civara",
    "analysis": {
      "scoreType": "average",
      "groupVariable": "gender",
      "groups": {
        "male": {
          "sampleSize": 482,
          "mean": 65.84,
          "standardDeviation": 14.10,
          "minimum": 9.00,
          "maximum": 100.00
        },
        "female": {
          "sampleSize": 518,
          "mean": 69.57,
          "standardDeviation": 14.54,
          "minimum": 23.00,
          "maximum": 100.00
        }
      },
      "comparison": {
        "meanDifference": -3.73,
        "standardError": 0.91,
        "degreesOfFreedom": 996.12,
        "confidenceLevel": 0.95,
        "confidenceInterval": {
          "lower": -5.52,
          "upper": -1.94
        },
        "tStatistic": -4.09,
        "pValue": 0.00005,
        "isSignificant": true
      },
      "interpretation": "The 95% confidence interval does not contain zero. Therefore, there is a statistically significant difference between the mean scores of male and female students at the 5% significance level."
    }
  },
  "meta": null
}
\end{lstlisting}

\section{Response Fields}

\begin{longtable}{p{5cm} p{8cm}}
\toprule
\textbf{Field} & \textbf{Description} \\
\midrule
\texttt{success} & Indicates whether the request was processed successfully. \\
\texttt{message} & A human-readable summary of the result. \\
\texttt{data} & The response payload. When \texttt{success=true}, this contains the expected result object or list. \\
\texttt{meta} & Optional metadata object. \texttt{null} when not applicable. \\
\texttt{meta.pagination} & Present only for list endpoints that support pagination. \\
\texttt{data.project} & The project name. \\
\texttt{data.analysis.scoreType} & The score used in the analysis. \\
\texttt{data.analysis.groupVariable} & The variable used to separate the two groups. \\
\texttt{sampleSize} & Number of observations in each group. \\
\texttt{mean} & Sample mean for each group. \\
\texttt{standardDeviation} & Sample standard deviation for each group. \\
\texttt{minimum} & Minimum score in each group. \\
\texttt{maximum} & Maximum score in each group. \\
\texttt{meanDifference} & Difference between the two group means. \\
\texttt{standardError} & Standard error of the difference between means. \\
\texttt{degreesOfFreedom} & Welch degrees of freedom. \\
\texttt{confidenceLevel} & Confidence level used in the interval. \\
\texttt{confidenceInterval.lower} & Lower bound of the confidence interval. \\
\texttt{confidenceInterval.upper} & Upper bound of the confidence interval. \\
\texttt{tStatistic} & Test statistic from Welch's two-sample t-test. \\
\texttt{pValue} & Probability value from the hypothesis test. \\
\texttt{isSignificant} & Indicates whether the result is statistically significant. \\
\texttt{interpretation} & Plain-language statistical conclusion. \\
\bottomrule
\caption{Analysis response fields}
\end{longtable}

\section{Statistical Method}

The API compares two independent groups using Welch's two-sample method.

The parameter of interest is:

\[
\mu_{\text{male}} - \mu_{\text{female}}
\]

The standard error is calculated as:

\[
SE =
\sqrt{
\frac{s_m^2}{n_m}
+
\frac{s_f^2}{n_f}
}
\]

The 95\% confidence interval is calculated as:

\[
(\bar{x}_m - \bar{x}_f)
\pm
t_{\alpha/2} \times SE
\]

where:

\begin{itemize}
    \item $\bar{x}_m$ is the male sample mean.
    \item $\bar{x}_f$ is the female sample mean.
    \item $s_m$ is the male sample standard deviation.
    \item $s_f$ is the female sample standard deviation.
    \item $n_m$ is the male sample size.
    \item $n_f$ is the female sample size.
\end{itemize}

The null and alternative hypotheses are:

\[
H_0: \mu_{\text{male}} = \mu_{\text{female}}
\]

\[
H_1: \mu_{\text{male}} \neq \mu_{\text{female}}
\]

The result is considered statistically significant when:

\[
p < 0.05
\]

\section{Validation Rules}

The API validates the uploaded dataset before analysis.

\begin{table}[h!]
\centering
\renewcommand{\arraystretch}{1.4}
\begin{tabular}{p{5cm} p{8cm}}
\toprule
\textbf{Validation Rule} & \textbf{Expected Behavior} \\
\midrule
Missing file & Return a validation error. \\
Invalid file type & Return a validation error. \\
Missing required columns & Return a validation error. \\
Non-numeric score values & Return a validation error. \\
Less than two gender groups & Return a validation error. \\
Empty dataset & Return a validation error. \\
\bottomrule
\end{tabular}
\caption{Dataset validation rules}
\end{table}

\section{Error Response Format}

All error responses follow the standard response envelope.

\begin{lstlisting}
{
  "success": false,
  "message": "Validation failed",
  "data": null,
  "meta": null,
  "error": {
    "code": "VALIDATION_ERROR",
    "message": "The uploaded file is missing the required column: gender."
  }
}
\end{lstlisting}

\section{Error Codes}

\begin{table}[h!]
\centering
\renewcommand{\arraystretch}{1.4}
\begin{tabular}{p{4cm} p{3cm} p{6cm}}
\toprule
\textbf{Error Code} & \textbf{HTTP Status} & \textbf{Description} \\
\midrule
\texttt{NO\_FILE} & 400 & No file was uploaded. \\
\texttt{\detokenize{INVALID_FILE_TYPE}} & 400 & The uploaded file is not a CSV file. \\
\texttt{\detokenize{MISSING_COLUMNS}} & 400 & One or more required columns are missing. \\
\texttt{\detokenize{INVALID_DATA}} & 400 & The dataset contains invalid or non-numeric values. \\
\texttt{\detokenize{INSUFFICIENT_GROUPS}} & 400 & The dataset does not contain two valid groups. \\
\texttt{\detokenize{ANALYSIS_FAILED}} & 500 & The server failed to complete the analysis. \\
\bottomrule
\end{tabular}
\caption{API error codes}
\end{table}

\section{Example Request}

Using \texttt{curl}:

\begin{lstlisting}
curl -X POST https://civara.mahros.dev/api/v1/analyze \
  -F "file=@students.csv" \
  -F "scoreType=average"
\end{lstlisting}

\section{Client-Side Cache Recommendation}

Since the API does not use a database, the frontend may store the latest result in browser storage.

Recommended browser storage key:

\begin{lstlisting}
civara:last-analysis
\end{lstlisting}

The cached object may contain:

\begin{lstlisting}
{
  "uploadedAt": "2026-05-12T10:30:00.000Z",
  "fileName": "students.csv",
  "analysis": {}
}
\end{lstlisting}

This allows the dashboard to reload the latest analysis without requiring the user to upload the dataset again.

\section{Conclusion}

This API contract defines a minimal and practical public interface for Civara. The design is intentionally simple: a user uploads a CSV dataset, the API validates and analyzes the data, and the response contains the statistical results needed for interpretation.

The API supports the main goal of the project by converting a manual statistical workflow into a small, usable software system.
\end{document}
